\documentclass[11pt,a4paper]{article}
\usepackage[margin=23mm]{geometry}
\usepackage{amsmath,amssymb,array,booktabs}
\input{glyphtounicode}
\pdfgentounicode=1
\usepackage[hidelinks]{hyperref}
\hypersetup{hypertexnames=false}
\setlength{\parindent}{0pt}
\setlength{\parskip}{5pt}
\newcommand{\pos}[1]{[#1]_+}
\renewcommand{\labelitemi}{\(\bullet\)}
\begin{document}
\begin{center}
{\Large S-E: local admission, causal transitions and an allocation gate}\\[4pt]
15 October 2026 --- analytical qualification, no traffic experiment
\end{center}

\textbf{Status.} S-E is the user-selected first qualification direction.
The common trial constitutive law and the S-E-P/S-E-W allocations below
are proposals, not approved task equations. Stop at \textbf{S-E-A2}, the
free-A allocation interpretation. The whole three-part review and Step 3
remain unfinished. No source theorem is transferred to this node.

\section*{1. Retained interface and a concrete trial state-to-work law}

H1 supplies terminal entrance histories $F_j$ with the accepted common
constant advice $v\in[v_-,v_+]$. The terminal has length $b$, initially
empty histories, and no internal lateral source. Put
\begin{equation}
 C(v)=\frac{wk_jv}{w+v},\quad
 P_j(t)=F_j(t-b/v),\quad q_j=P_j-Y_j,\quad N_j=F_j-Y_j,
 \quad j=A,E,B. \tag{1}
\end{equation}
Let $\lambda_j=\dot P_j$ denote eligible arrival rates. B retains
$Y_B=P_B$ on its independent adequate-service domain. For A/E, let $B(t)$
denote available shared work, $0\leq B\leq R\leq C(v_-)$.
Here $B$ is an input to an instantaneous audit, not an externally supplied
future receiving trace. Its eventual law must be autonomous and causal.

Use mode $m_j=F$ for free/unqueued admission and $m_j=Q$ for a selected
queued boundary trace. Initially $q_j=0,m_j=F$. An active queue has demand
$C(v)$; an empty free approach offers $\lambda_j$. The proposed equilibrium
trace identification is
\begin{equation}
 u_j=v\ (m_j=F),\qquad
 u_j=U(g_j):=\frac{wg_j}{wk_j-g_j}\ (m_j=Q),\quad
 0\leq g_j\leq C(v),\quad U(C(v))=v. \tag{2}
\end{equation}
An absent free stream's formal $v$ is not a measured speed. Equation (2)
is an upstream boundary trace, not mean lane speed. Identifying that trace
with a representative insertion/recovery mark remains an adaptation.

One concrete trial constitutive family is
\begin{equation}
 c_A=1,\qquad
 c_E=1+\delta+\beta\frac{g_A}{C(v)}
          \frac{\pos{u_A-u_E}}v,\qquad \delta,\beta\geq0,
 \qquad g_A+c_Eg_E\leq B. \tag{3}
\end{equation}
$c_j$ is work per admitted mass. $\delta$ represents immediate maneuver
work; it may be zero. The additional one-sided term represents harder
entry when A is faster, weighted by A's actual passing flow. No absent A
stream produces a mismatch penalty. This is not an empirically identified
formula or a universal sign theorem. The finite late work and its separation
from post-entry recovery need physical qualification. $c_A=1$ normalizes
ordinary current work; it does not forbid a justified A recovery influence.

\newpage
\section*{2. What is unique inside a selected regime}

Retain the proposed fair-work allocator as a \emph{comparison template}:
for demands $d_A,d_E$, costs $1,c_E$ and budget $B$, choose minimal $z\geq0$
such that
\begin{equation}
 \min(d_A,z)+\min(c_Ed_E,z)
   =\min(B,d_A+c_Ed_E),\quad
 g_A=\min(d_A,z),\quad g_E=\min(d_E,z/c_E). \tag{4}
\end{equation}
This is equal work with unused-work redistribution, not Carey's
demand-proportional R2 or a literature-derived fairness theorem.

When E is free, $u_E=v$ and $c_E=c_0:=1+\delta$ for either A trace.
Thus (4) has the previously proved unique fixed-input allocation.
When E is queued, its work offer reaches $B$. Under (4),
\begin{equation}
 a:=g_A=\min(d_A,B/2),\qquad x:=B-a,\qquad
 H_E(g):=g\!\left[c_0+\beta\frac a{C(v)}
                    \frac{\pos{u_A-U(g)}}v\right]=x. \tag{5}
\end{equation}
For a free-A-priority branch with A still served freely, instead take
$a=\lambda_A\leq B$ and $x=B-a$ in (5).
In both cases $u_A$ is fixed by that selected A mode and $a$.

A conservative sufficient condition for a unique feasible root is
\begin{equation}
 \eta:=1+\delta-\beta(1+v_+/w)>0. \tag{6}
\end{equation}
Indeed $U'(g)=w^2k_j/(wk_j-g)^2$ and
$gU'(g)\leq v(1+v/w)$ on $[0,C(v)]$. On the active positive-part branch,
\begin{equation}
 H_E'(g)=c_0+\beta\frac a{C(v)v}
                  \{u_A-U(g)-gU'(g)\}\ \geq\eta; \tag{7}
\end{equation}
off it $H_E'=c_0$. At the kink the continuous function remains strictly
increasing. Also $H_E(0)=0$ and $H_E(C(v))=c_0C(v)\geq B\geq x$.
Consequently (5) has exactly one root in $[0,C(v)]$, with
\begin{equation}
 \frac{x}{c_0+\beta}\leq g_E\leq\frac{x}{c_0},\qquad
 0\leq g_A+g_E\leq B. \tag{8}
\end{equation}
Condition (6) is a mathematical sufficient domain, not an approved
parameter restriction on real traffic. It need not be necessary.
No parameter is fitted or chosen for a positive control result.

\textbf{Limit check.} With both streams queued under (4), $a=B/2$ and
$g_E\leq a$. If $\delta=0$, the unique solution is $g_E=a$ with matched
traces and $c_E=1$. If traces match with $\delta>0$, the mismatch term
vanishes but the declared maneuver work remains. Recovery may still reduce
$B$. Comparable speeds do not by themselves prove high total service.

\newpage
\section*{3. Two allocations and the causal transition rules}

\textbf{S-E-W, immediate cooperation.} Apply (4) in every mode. Free A
does not have an additional protection rule. Its service can be reduced
while all E arrivals pass. Both queued streams retain positive service
when $B>0$.

\textbf{S-E-P, free-A priority then queued cooperation.} While A is empty
and its free offer fits $B$, put
\begin{equation}
 g_A=\lambda_A,\quad
 g_E=\begin{cases}
 \min\{\lambda_E,(B-\lambda_A)/c_0\},&m_E=F,\\
 H_E^{-1}(B-\lambda_A),&m_E=Q.
 \end{cases} \tag{9}
\end{equation}
If A cannot fit $B$, its strict shortfall triggers its queued mode.
Whenever A is queued, use (4)--(5), including a free E offer if applicable.
This proposed regime adaptation is not strict A priority in all states:
with both queues nonempty, A does not consume all available service.

For either allocation the proposed S-E transition semantics are:
\begin{enumerate}
\item Test admission in the current free modes. If every empty approach's
eligible offer is met, keep those approaches free. An alternative queued
root is not an activation condition.
\item A strict free-mode shortfall selects the affected queued trace;
recompute the coupled node and require nonnegative right-hand queue growth
for a newly active empty queue. In the strict onset examples below, this
condition holds. No external seed or numerical queue epsilon is used.
\item While $q_j>0$, evolve $\dot q_j=\lambda_j-g_j$ and
$\dot Y_j=g_j$. At a true zero hit, test free admission again. Never serve
unavailable mass. Equality alone does not trigger a new queued branch.
\item Retain the common recovery history across queue clearance. At load
jumps use physical right-hand data. Simultaneous events and grazing need
qualification in the complete coupled law.
\end{enumerate}

These rules reject the old arbitrary-birth diagnostic: its free work demand
is $7/10<4/5$, so no activation is available under the constant local inputs.
That is a local exclusion of the specific delayed births, not a global
well-posedness theorem. The fixed-regime inverse bound (6) and strict-event
checks do not prove uniqueness through every coupled switching event.

The distinction between (9) and (4) selects which stream first waits and
which departure speeds feed later recovery. S-E direction approval does
not select that allocation, an exact equal-work convention, or (3).

\newpage
\section*{4. A reachable H1 front exposes the allocation difference}

Reuse the accepted joined H1/T1 algebraic family, in units with its reference
speed and density equal to one:
\begin{equation}
 \begin{gathered}
 v\in[1,2],\quad w=2,\quad k_j=3,\quad
 e_0=7/10,\quad a_0=b_0=1/4,\quad\gamma=2,\\
 R=39/20,\quad R_B=1,\quad b=\ell>0,\quad
 L_p=\ell+h\log2,\quad L=L_p+b,\quad h>0.
 \end{gathered} \tag{10}
\end{equation}
The existing exact H1 preparation gives front ending density
$e_f(2)=125/352$. At the first eligible front for $v=2$,
\begin{equation}
 C=3,\quad \lambda_E=125/176,\quad
 \lambda_A=1047/880,\quad\lambda_A+\lambda_E=19/10<R. \tag{11}
\end{equation}
The terminal and history are empty, hence $B=R$ for an initially neutral
continuous recovery account. For this exact static comparison only, take
$\delta=\beta=1/10$ in (3); these values are not selected task parameters.
Then $\eta=9/10>0$ and
\begin{equation}
 c_0\lambda_E=25/32<R/2,\qquad
 \lambda_A+c_0\lambda_E-R=37/1760>0. \tag{12}
\end{equation}

\textbf{S-E-W.} All E arrivals fit their fair-work share. The causal onset
is A waiting, with
\begin{equation}
 g_E=125/176,\quad g_A=187/160,\qquad
 \dot q_A=37/1760>0,\quad\dot q_E=0. \tag{13}
\end{equation}
The queued A trace becomes $374/773$, while free E retains $v=2$ and
$c_E=c_0$. This is a compatible right-hand regime for the trial law.

\textbf{S-E-P.} Free A first passes its full offer. The free-E test gives
\begin{equation}
 g_E^{F}=669/968<125/176,\qquad
 \lambda_E-g_E^{F}=37/1936>0. \tag{14}
\end{equation}
Thus the declared maneuver-work shortfall activates queued E. Its rate solves
\begin{equation}
 g\left[\frac{11}{10}+\frac1{10}\frac{1047/880}{3}
                  \left(1-\frac{g}{6-g}\right)\right]=\frac{669}{880},
 \quad 5189g^2-33432g+20070=0. \tag{15}
\end{equation}
Only
\begin{equation}
 g_* =\frac{16716-183\sqrt{5234}}{5189},\quad
 \frac{2230}{3343}<g_*<\frac{669}{968}<\frac{125}{176} \tag{16}
\end{equation}
is feasible. A remains free; E has a strictly positive queue-growth margin.
The strict inequalities persist locally under continuous histories and
right-continuous arrival data. This is an H1-reachable local onset check,
not a time-stepped traffic trajectory or complete cohort result.

\newpage
\section*{5. A simpler demand-proportional graft fails the mode test}

Carey's R2 and demand-based merge models are important simpler rivals.
Their native discretization and a properly invariant LTM junction cannot
be identified with the following naive graft. Test, only for this audit,
\begin{equation}
 g_A=\sigma d_A,\quad g_E=\sigma d_E,\quad
 \sigma=\min\!\left\{1,\frac{B}{d_A+c_Ed_E}\right\},\quad
 d_j=\begin{cases}\lambda_j,&m_j=F,\\ C,&m_j=Q.\end{cases} \tag{B1}
\end{equation}
At the same initially empty H1 front (11), there is no compatible strict
right-hand mode under these binary demand/trace rules:
\begin{itemize}
\item Both free: $\sigma=3432/3469<1$. Both arrivals are restricted,
so neither empty approach can stay free with zero waiting derivative.
\item Only A queued: E is free, so $c_E=c_0$ and
$g_A=936/605>1047/880=\lambda_A$. A's newly active empty queue would
have a negative derivative. E's free offer is also restricted.
\item Only E queued: $1.1\leq c_E\leq1.2$ gives
\begin{equation}
 g_E=\frac{3B}{\lambda_A+3c_E}
       \geq\frac{1716}{1405}>\frac{125}{176}=\lambda_E. \tag{B2}
\end{equation}
The newly active empty E queue would have a negative derivative.
\item Both queued: equal mass demands give equal actual flows and hence
equal congested traces. Thus $c_E=c_0$ and
\begin{equation}
 g_A=g_E=\frac{B}{1+c_0}=\frac{13}{14}
                   >\frac{125}{176}=\lambda_E. \tag{B3}
\end{equation}
Again E's newborn queue cannot have that negative derivative.
\end{itemize}
The change from an empty demand to the capacity offer immediately alters
the proportional allocation. Merely saying ``activate the queue on a
shortfall'' does not complete (B1). A numerical method might alternate
between branches, but its step must not select a scientific flux law.

This excludes the specific binary-demand graft, not all demand-based
merging, native R2 or every possible generalized solution. An explicit
interior-state/selection law would be an additional model. The primary
LTM source itself distinguishes invariant and non-invariant fair merges
and gives a stationary counterexample for the latter in section 6.
Our transient work example is a separate deduction; its proof is not
attributed to that paper. Neither S-E-P nor S-E-W is declared generally
invariant or globally well posed on the strength of this test.

\newpage
\section*{6. Joined implications, proof limits and the stopping gate}

\textbf{Physical interface.} In (10), $M=6\ell/5<3\ell=k_jb$ and
all H1 entrance rates fit $C(v)$ across $[1,2]$. For any eventual admissible
departures, the retained finite-terminal vacancy bound gives
$V_j(t)\geq k_jb-M>0$ after $b/w$; the early-time bound follows as before
from $\dot F_j\leq C(v)<wk_j$. The handoff protection does not depend
on which allocation is selected. No point-buffer substitution is made.
Conservation and complete local TST retain
\begin{equation}
 N_j=F_j-Y_j,\qquad
 J(v)=\frac{M(L-\ell/2)}v+\int_0^\infty(q_A+q_E+q_B)dt. \tag{17}
\end{equation}
Completed mass leaves once; service memory is not physical stock.

\textbf{Recovery is not closed by the admission check.} The two onsets
place slow boundary traces in different streams. E-only recovery would
be neutral at the initial S-E-W onset; that does not prove its whole shared
account neutral. A's actual slow passage may have a justified continuing
influence. Neither deleting that influence nor charging every queued A
departure independently is permitted without qualification. Both allocations
retain one account, not independent ceilings per past departure.

For orientation, the earlier isolated kernel may be marked by a justified
actual insertion speed $u$, rather than a fixed E speed. If only a separately
justified E source were present, a candidate $D(\theta;u)$ would satisfy
\begin{equation}
 p_E(t)=m_{\rm ref}^{-1}\!\int_0^tD(t-s;u_E(s))dY_E(s),\quad
 0\leq D(\theta;u)\leq D(\theta;0),\quad
 \int_0^\infty p_E(t)dt\leq\frac{MW(0)}{m_{\rm ref}}. \tag{18}
\end{equation}
This is a conditional extension bound, not an adopted E-only meter,
overlap rule or full M2 proof. If an admitted full account establishes
$\int(R-B)dt\leq\Lambda$ and $c_E\leq c_{\max}=1+\delta+\beta$,
either work-conserving allocation clears after last eligibility $t_*=L/v$
by the previously derived sufficient bound
\begin{equation}
 t_*+\frac{c_{\max}(q_A(t_*)+q_E(t_*))+\Lambda}{R}. \tag{19}
\end{equation}
No positive service floor or transferred M1 uniqueness proof is used.

\textbf{Evidence disposition.} Proved: unique fixed-regime roots under (6),
specific causal-onset consistency, rejection of the old spontaneous-birth
diagnostic, and a reachable local H1 allocation contrast. Retained:
conditional storage/accounting/clearance interfaces. Not proved: whole-law
uniqueness through simultaneous/grazing events, the actual joint recovery
law, complete-cohort MP6, observed merge credibility, an action split or
gradient usefulness. Conditional scalar reconstruction remains a strong rival.

\textbf{Gate S-E-A2 (Type 2 G3).} Choose whether free A first claims its
offer until reduced shared service makes it queue (S-E-P), or cooperation
may slow A from the first conflict (S-E-W). Provisional first qualification:
S-E-P, retaining S-E-W as a serious comparator. Neither allocation is selected.
This
does not reopen S-E versus S-V, H1, the primary demand domain or the local
objective. The exact trial cost/trace identification remains a hypothesis.
Stop before selecting the allocation and dependent shared recovery.
\end{document}
